\documentclass[10pt,a4paper]{amsart}
\usepackage[margin=19mm]{geometry}
\usepackage{amsmath,amssymb,mathtools}
\usepackage[scr=rsfs,cal=euler]{mathalfa}
\usepackage{xfrac}
\usepackage[unicode,hidelinks,bookmarks=false]{hyperref}
\newcommand{\ie}{i.e.\ }
\newcommand{\cf}{cf.\ }
\newcommand{\resp}{respectively\ }
\newcommand{\Subsection}{\subsection}
\providecommand{\nobreakdash}{\nobreak\hbox{-}}
\setlength{\emergencystretch}{2em}
\multlinegap0pt
\usepackage{amssymb}
\usepackage{cancel}
\usepackage[all]{xy}
\usepackage{xcolor}
\usepackage{tikz}
\usepackage{csquotes}
\usepackage{bbm}
\usepackage{mathtools}
\usepackage{caption}
\usepackage{booktabs}
\newtheorem{lem}{}
\newtheorem{thm}{}
\title{Combinatorics behind discriminants of polynomial systems}
\author{Vladislav Pokidkin}
\date{Source: arXiv:2509.02963v2 (CC BY 4.0)}
\begin{document}
\maketitle

\noindent\emph{Source and licence.} Combinatorics behind discriminants of polynomial systems. Version arXiv:2509.02963v2 (CC BY 4.0), distributed by arXiv under the Creative Commons Attribution 4.0 International licence, http://creativecommons.org/licenses/by/4.0/, as recorded in the arXiv licence metadata for this exact version and shown on the abstract page https://arxiv.org/abs/2509.02963v2. Full licence notice: \texttt{licenses/arxiv-2509.02963-CC-BY-4.0.txt}.\par\smallskip

\noindent\textit{Editorial scope.} Counted unit: the article's thm BK (source0.tex lines 479--482). The article's complete original proof of it is delivered with it (source0.tex lines 484--505, one \texttt{proof} environment of 22 lines). No mathematics is written, repaired or re-derived: the statement and the proof are the article's own lines, in the article's own notation, and no retained line is substituted or re-flowed.  Retained uncounted as the source-local context the proof needs: lines 311--315; lines 390--394; lines 470--473.  Nothing was omitted from any retained span, so the package carries no omission record.  No retained line contains a citation, so the wrapper carries no bibliography and no bibliography entry.  No retained line contains a figure environment, an image-inclusion command or drawing code, so no image is delivered and the package declares no asset dependency.  Editorial formatting only, in addition: an AMS \texttt{amsart} preamble that restates the article's own declarations and macros used by the retained lines, namely the environments \texttt{lem}, \texttt{thm} on separate counters, together with the standard packages those lines need.  The retained environments print in this document as Lemma 1, Theorem 1 in order of appearance, whereas the article prints the same items with the numbering of its own full document.  The complete native source of the article as distributed in the arXiv e-print for this exact version is bundled unchanged as \texttt{sources/184381/source0.tex}.
\begin{lem}
\label{Lemma. C=00005Ce is a BK-set}For an element $e$ from a circuit
$C$ of the induced matroid from a polymatroid, the complement $C\backslash e$
is a BK-set.
\end{lem}

\begin{lem}
\label{Lemma. The union and intersection of BK-sets are BK-sets}In
an independent polymatroid, unions and intersections of BK-sets are
BK-sets.
\end{lem}

\begin{lem}
\label{Lemma. Cyclics are unions of BK-sets}In the induced matroid,
cycles are unions of BK-sets.
\end{lem}

\begin{thm}
\label{Theorem. In cyclic, bases are BK}In a cyclic polymatroid,
bases are BK-sets.
\end{thm}

\begin{proof}
A cyclic ground set $E$ is a union of BK-sets by Lemma \ref{Lemma. Cyclics are unions of BK-sets}
and is a union of circuits by definition. For a fixed element $e\in E,$
we can represent the ground set $E$ via the union of circuits as
follows:
\[
E=\left(\underset{e\in C}{\cup}C\right)\cup\left(\underset{e\notin C'}{\cup}C'\right).
\]
The union $\underset{e\notin C'}{\cup}C'$ is a cycle; hence, it is
a union of BK-sets by Lemma \ref{Lemma. Cyclics are unions of BK-sets}.
According to Lemma \ref{Lemma. C=00005Ce is a BK-set}, $C\backslash e$
is a BK-set for every circuit $C$ containing $e.$ Therefore, the
complement
\[
E\backslash e=\left(\underset{e\in C}{\cup}(C\backslash e)\right)\cup\left(\underset{e\notin C'}{\cup}C'\right)
\]
is a union of BK-sets. If the complement $E\backslash e$ is dependent,
we repeat the procedure and remove one by one elements until we obtain
an independent set $I\subset E$. This independent set $I$ is a basis
automatically, is a union of BK-sets by induction, and is a BK-set
by Lemma \ref{Lemma. The union and intersection of BK-sets are BK-sets}.
\end{proof}

\end{document}
